% ============================================================
% NDH Lineage Adjacency
% Module: 07-lineage-adjacency.tex
% Altitude: A4–A7 (Governance → Lineage → Identity)
% Membrane: Lineage-Constitution • Reversible Adjacency
% ============================================================

\section{Lineage Adjacency}
\label{sec:lineage-adjacency}

The Lineage Constitution governs reversible adjacency between the NDH
manifold and external systems. It defines lineage channels, reversible
interface morphisms, and non-absorptive lineage operators. Lineage adjacency
ensures that NDH may interact with external structures without inheriting,
absorbing, or collapsing identity.

% ------------------------------------------------------------
% 1. Lineage Channels
% ------------------------------------------------------------

\subsection{Lineage Channels}

A lineage channel is defined as a pair of functors

\[
\Lambda = (\lambda_{\mathrm{out}}, \lambda_{\mathrm{in}})
\]

where

\[
\lambda_{\mathrm{out}} : \mathcal{M} \to \mathcal{M}', \qquad
\lambda_{\mathrm{in}}  : \mathcal{M}' \to \mathcal{M}.
\]

The external manifold $\mathcal{M}'$ may represent any constitutionally
permitted external system.

Lineage channels must be reversible:

\[
\lambda_{\mathrm{in}} \circ \lambda_{\mathrm{out}} = \mathrm{id}_{\mathcal{M}}.
\]

Reversibility prohibits identity inheritance.

% ------------------------------------------------------------
% 2. Reversible Interface Morphisms
% ------------------------------------------------------------

\subsection{Reversible Interface Morphisms}

For any object $x \in \mathcal{M}$, the lineage interface morphism is

\[
\iota_x : x \leftrightarrow \lambda_{\mathrm{out}}(x).
\]

Interface morphisms must satisfy:

\begin{itemize}
    \item \textbf{Identity preservation}:
    \[
    \pi(\iota_x) = \mathrm{id}_{\mathcal{M}}.
    \]

    \item \textbf{Non-absorptive behavior}:
    \[
    \iota_x(a_1) = \iota_x(a_2) \Rightarrow a_1 = a_2.
    \]

    \item \textbf{Reversibility}:
    \[
    \lambda_{\mathrm{in}}(\lambda_{\mathrm{out}}(x)) = x.
    \]
\end{itemize}

These constraints ensure lineage safety.

% ------------------------------------------------------------
% 3. Lineage Adjacency Category
% ------------------------------------------------------------

\subsection{Lineage Adjacency Category}

Lineage adjacency is defined as the category

\[
\mathcal{A}_{\mathrm{lin}} = (\mathrm{Obj}_{\mathrm{lin}},
                              \mathrm{Hom}_{\mathrm{lin}})
\]

where objects are lineage states and morphisms represent reversible lineage
transitions.

Morphisms must satisfy:

\[
f \in \mathrm{Hom}_{\mathrm{lin}} \Rightarrow
\lambda_{\mathrm{in}} \circ f \circ \lambda_{\mathrm{out}}
= \mathrm{id}.
\]

This ensures that lineage transitions do not accumulate or absorb identity.

% ------------------------------------------------------------
% 4. Lineage Operators
% ------------------------------------------------------------

\subsection{Lineage Operators}

Lineage operators act on lineage adjacency:

\[
\mathcal{O}_{\mathrm{lin}} :
\mathcal{A}_{\mathrm{lin}} \to \mathcal{A}_{\mathrm{lin}}.
\]

They must satisfy:

\begin{itemize}
    \item \textbf{Reversibility preservation}:
    \[
    \mathcal{O}_{\mathrm{lin}}(f) \text{ is reversible}.
    \]

    \item \textbf{Identity preservation}:
    \[
    \pi \circ \mathcal{O}_{\mathrm{lin}} = \pi.
    \]

    \item \textbf{Non-absorptive lineage flow}:
    \[
    \mathcal{O}_{\mathrm{lin}}(f_1) = \mathcal{O}_{\mathrm{lin}}(f_2)
    \Rightarrow f_1 = f_2.
    \]
\end{itemize}

Lineage operators enforce constitutional lineage safety under $\mathcal{C}_6$.

% ------------------------------------------------------------
% 5. Interaction with Adjacency Categories
% ------------------------------------------------------------

\subsection{Adjacency Interaction}

Lineage adjacency interacts with NDH adjacency categories via lineage lifts:

\[
\mathcal{L}_s : \mathcal{A}_s \to \mathcal{A}_{\mathrm{lin}}, \qquad
\mathcal{L}_c : \mathcal{A}_c \to \mathcal{A}_{\mathrm{lin}}, \qquad
\mathcal{L}_o : \mathcal{A}_o \to \mathcal{A}_{\mathrm{lin}}.
\]

These lifts ensure that spectral, constitutional, and operator-level
adjacencies respect reversible lineage constraints.

% ------------------------------------------------------------
% 6. Interaction with Quantum Sponge
% ------------------------------------------------------------

\subsection{Quantum Sponge Interaction}

The Quantum Sponge interacts with lineage adjacency via the functorial lift

\[
\mathcal{Q}_{\mathrm{lin}} :
\mathcal{A}_{\mathrm{lin}} \to \mathbf{Cat}.
\]

This ensures that reconstruction behavior remains lineage-safe and
non-absorptive.

% ------------------------------------------------------------
% 7. Constitutional Integration
% ------------------------------------------------------------

\subsection{Constitutional Integration}

The Lineage Constitution is represented by the functor $\mathcal{C}_6$:

\[
\mathcal{C}_6 : \mathcal{M} \to \mathcal{M}.
\]

It must satisfy:

\[
\mathcal{C}_6(\lambda_{\mathrm{out}}) = \lambda_{\mathrm{out}}, \qquad
\mathcal{C}_6(\lambda_{\mathrm{in}})  = \lambda_{\mathrm{in}}.
\]

This integrates lineage adjacency with the constitutional functor tower.

% ------------------------------------------------------------
% End of Module
% ------------------------------------------------------------
